Considereu les dues matrius següents:
\[
A=\begin{pmatrix}2&-3&-5\\-1&4&5\\1&-3&-4\end{pmatrix},\qquad B=\begin{pmatrix}2&2&0\\-1&-1&0\\1&2&1\end{pmatrix}.
\]

\begin{apartats}

\apartat{1,5}
Calculeu les matrius $A\cdot B$ i $B\cdot A$.

\begin{solucio}
\[
A\cdot B=\begin{pmatrix}2&-3&-5\\-1&4&5\\1&-3&-4\end{pmatrix}\cdot\begin{pmatrix}2&2&0\\-1&-1&0\\1&2&1\end{pmatrix}
=\begin{pmatrix}2&-3&-5\\-1&4&5\\1&-3&-4\end{pmatrix}=A,
\]
\[
B\cdot A=\begin{pmatrix}2&2&0\\-1&-1&0\\1&2&1\end{pmatrix}\cdot\begin{pmatrix}2&-3&-5\\-1&4&5\\1&-3&-4\end{pmatrix}
=\begin{pmatrix}2&2&0\\-1&-1&0\\1&2&1\end{pmatrix}=B .
\]
La resolució es donarà per bona encara que no apareguin els càlculs intermedis.

\textit{Pauta oficial:} 0,75 pel càlcul del producte $A\cdot B$ i 0,75 pel càlcul del producte $B\cdot A$.
\end{solucio}

\apartat{1}
Siguin $C$ i $D$ dues matrius quadrades del mateix ordre que satisfan $C\cdot D=C$ i $D\cdot C=D$. Comproveu que les dues
matrius, $C$ i $D$, són idempotents.

\emph{Nota:} Una matriu quadrada s'anomena \emph{idempotent} si coincideix amb el seu quadrat.

\begin{solucio}
Sabem que $(*)$ $C\cdot D=C$ i $(**)$ $D\cdot C=D$. Aleshores,
\[
C^2=C\cdot C\underset{(*)}{=}(C\cdot D)\cdot C=C\cdot(D\cdot C)\underset{(**)}{=}C\cdot D\underset{(*)}{=}C,\qquad
D^2=D\cdot D\underset{(**)}{=}(D\cdot C)\cdot D=D\cdot(C\cdot D)\underset{(*)}{=}D\cdot C\underset{(**)}{=}D .
\]
I, per tant, efectivament les dues matrius, $C$ i $D$, són idempotents.

\textit{Pauta oficial:} 0,5 per la comprovació que la matriu $C$ és idempotent i 0,5 per la comprovació que la matriu
$D$ és idempotent.
\end{solucio}

\end{apartats}
